\section{Introducción: todo es real}

La conversación comienza con una apertura evocadora: Ilya expresa asombro de que "todo esto sea real", refiriéndose a todo lo que está sucediendo en el campo de la IA. Dwarkesh añade una observación igualmente profunda: \textbf{qué tan normal se siente el slow takeoff}: el 1\% del PIB mundial se está invirtiendo en IA, pero esto solo se manifiesta en las noticias como "tal empresa anunció un monto de inversión difícil de comprender"; la persona común no ha sentido nada sustancial.

\begin{knowledgebox}{El efecto psicológico del Slow Takeoff}
Ilya considera que el impacto de la IA \textbf{se sentirá}: la IA se difundirá a través de la economía, con poderosas fuerzas económicas impulsando este proceso. Pero Dwarkesh plantea: "incluso al entrar a la singularity (singularity), la perspectiva de la persona común quizá no sea muy distinta." Ilya no está de acuerdo con esto: considera que el impacto terminará por percibirse profundamente.
\end{knowledgebox}

